\documentclass[10pt,a4paper]{amsart}
\usepackage[margin=19mm]{geometry}
\usepackage{amsmath,amssymb,mathtools}
\usepackage[scr=rsfs,cal=euler]{mathalfa}
\usepackage{xfrac}
\usepackage[unicode,hidelinks,bookmarks=false]{hyperref}
\newcommand{\ie}{i.e.\ }
\newcommand{\cf}{cf.\ }
\newcommand{\resp}{respectively\ }
\newcommand{\Subsection}{\subsection}
\providecommand{\nobreakdash}{\nobreak\hbox{-}}
\setlength{\emergencystretch}{2em}
\multlinegap0pt
\usepackage{mdframed}
\usepackage{mdframed}
\usepackage{mdframed}
\usepackage{mdframed}
\usepackage{xcolor}
\usepackage{mathtools}
\usepackage{amssymb}
\usepackage{enumerate}
\usepackage[normalem]{ulem}
\usepackage{csquotes}
\usepackage{mdframed}
\newtheorem{lemma}{Lemma}
\newcommand{\VV}{\mathbb{V}} % integers
\newcommand{\abs}[1]{\left| #1 \right|}
\renewcommand{\emptyset}{\varnothing} % emptyset -> varnothing
\newcommand{\set}[1]{\left\{ #1 \right\}}
\newcommand{\tup}[1]{\left( #1 \right)}
\title{Degree Sequences vs. Forests in Finite Graphs}
\author{Benjamin Liber}
\date{Source: arXiv:2603.02151v2 (CC BY 4.0)}
\begin{document}
\maketitle

\noindent\emph{Source and licence.} Degree Sequences vs. Forests in Finite Graphs. Version arXiv:2603.02151v2 (CC BY 4.0), distributed by arXiv under the Creative Commons Attribution 4.0 International licence, http://creativecommons.org/licenses/by/4.0/, as recorded in the arXiv licence metadata for this exact version and shown on the abstract page https://arxiv.org/abs/2603.02151v2. Full licence notice: \texttt{licenses/arxiv-2603.02151-CC-BY-4.0.txt}.\par\smallskip

\noindent\textit{Editorial scope.} Counted unit: the article's lemma vector-list (source0.tex lines 771--775). The article's complete original proof of it is delivered with it (source0.tex lines 777--860, one \texttt{proof} environment of 84 lines). No mathematics is written, repaired or re-derived: the statement and the proof are the article's own lines, in the article's own notation, and no retained line is substituted or re-flowed.  No further source-local context is needed: every cross-reference of every retained line resolves inside the two delivered spans.  Nothing was omitted from any retained span, so the package carries no omission record.  No retained line contains a citation, so the wrapper carries no bibliography and no bibliography entry.  No retained line contains a figure environment, an image-inclusion command or drawing code, so no image is delivered and the package declares no asset dependency.  Editorial formatting only, in addition: an AMS \texttt{amsart} preamble that restates the article's own declarations and macros used by the retained lines, namely the environments \texttt{lemma} on separate counters, together with the standard packages those lines need.  The retained environments print in this document as Lemma 1 in order of appearance, whereas the article prints the same items with the numbering of its own full document.  The complete native source of the article as distributed in the arXiv e-print for this exact version is bundled unchanged as \texttt{sources/195602/source0.tex}.
\begin{lemma}\label{vector-list}
    For every finite list $X$ of real vectors, \begin{align}\label{subset_sums_index_sets}
        \abs{\Sigma(X)} \geq \abs{\mathcal{I}(X)}.
    \end{align}
\end{lemma}

\begin{proof}
    Let us induct on $m=\abs{X}$.
    \medskip
    
    \textit{Base case:} The case $m=0$ is immediate: In this case, $\Sigma(X) = \set{\mathbf{0}}$ and $\mathcal{I}\tup{X} = \set{\emptyset}$ (since the empty list of vectors is linearly independent), so that both sides of \eqref{subset_sums_index_sets} equal $1$. This completes the base case.
    % First, suppose $m=1$. Then $X = \tup{\mathbf{x}_1}$ and $[m] = [1] = \set{1}$. There are exactly two subsets of $[1]$: $\emptyset$ and $\set{1}$. Taking $\mathbf{0} \in \VV$ to denote the zero vector, we have \begin{align*}
        % \sum_{i \in \emptyset}\mathbf{x}_i = \mathbf{0} \quad \text{and} \quad \sum_{i \in \set{1}}\mathbf{x}_i = \mathbf{x}_1;
    % \end{align*} hence, $\Sigma(X) = \set{\mathbf{0},\mathbf{x}_1}$.

    % Now, consider $\mathcal{I}(X)$. By convention, $\set{\mathbf{x}_i}_{i \in \emptyset} = \emptyset$ is linearly independent, and $\set{\mathbf{x}_i}_{i \in \set{1}} = \set{\mathbf{x}_1}$ is linearly independent if and only if $\mathbf{x}_1 \neq \mathbf{0}$. So, if $\mathbf{x}_1 = \mathbf{0}$, then $\mathcal{I}(X) = \set{\emptyset}$; otherwise, $\mathcal{I}(X) = \set{\emptyset,\set{1}}$. In both cases, \begin{align*}
        % \abs{\Sigma(X)} = \abs{\mathcal{I}(X)},
    % \end{align*} so that \eqref{subset_sums_index_sets} holds for the base case.

    \medskip
    \textit{Induction step:} Next, suppose \eqref{subset_sums_index_sets} holds for all lists of $m-1$ vectors in a real vector space. Let $X = \tup{\mathbf{x}_1,\ldots,\mathbf{x}_{m-1},\mathbf{x}_m}$ be a list of $m$ vectors in $\VV$. Set \begin{align*}
        X' := \tup{\mathbf{x}_1,\ldots,\mathbf{x}_{m-1}}, \quad \mathbf{x} := \mathbf{x}_m,
    \end{align*} so that $X = \tup{X',\mathbf{x}}$ and $|X'|=m-1$.

    First, suppose $\mathbf{x} = \mathbf{0}$. Then \eqref{subset_sums_index_sets} follows immediately: adjoining the zero vector does not create any new subset sums, and no linearly independent subset of $X$ contains $\mathbf{0}$. Hence, \begin{align*}
        \Sigma(X) = \Sigma(X') \quad \text{and} \quad \mathcal{I}(X) = \mathcal{I}(X'),
    \end{align*} so the result follows from the induction hypothesis applied to $X'$.

    And so, we may assume that $\mathbf{x} \neq \mathbf{0}$. Let $S := \Sigma(X')$. By the induction hypothesis, \begin{align}\label{induct_X}
        \abs{S} = \abs{\Sigma(X')} \geq \abs{\mathcal{I}(X')}. 
\end{align} Moreover, \[\Sigma(X) = S \cup (S+\mathbf{x}),\] where $S+\mathbf{x} = \set{\mathbf{s}+\mathbf{x} : \mathbf{s} \in S}$. Indeed, every subset sum either excludes $\mathbf{x}$, yielding an element of $S$, or includes $\mathbf{x}$, yielding an element of $S+\mathbf{x}$. Hence, \begin{align}\label{split_sigmaX}
    \abs{\Sigma(X)} = \abs{S} + \abs{(S+\mathbf{x}) \setminus S}.
\end{align}

Next, consider the $1$-dimensional vector subspace \begin{align*}
    \langle \mathbf{x} \rangle = \set{c\mathbf{x} : c \in \mathbb{R}}
\end{align*} of $\VV$, and let $\pi : \VV \to \VV / \langle \mathbf{x} \rangle$ denote the quotient map, given by \begin{align*}
    \pi(\mathbf{v}) = \mathbf{v}+ \langle \mathbf{x} \rangle \quad \text{for all } \mathbf{v} \in \VV.
\end{align*} In particular, since $\mathbf{x} \in \langle \mathbf{x} \rangle$, we have $\pi(\mathbf{x}) = \mathbf{0}$. 

Let $\ell : \VV \to \mathbb{R}$ be a linear functional such that $\ell(\mathbf{x}) > 0$ (since $\mathbf{x} \neq \mathbf{0}$, such a linear functional exists). For each $y \in \pi(S)$, let \begin{align*}
    F_y := S \cap \pi^{-1}(y),
\end{align*} which is finite and nonempty. Choose $\mathbf{s}_y \in F_y$ such that \begin{align*}
    \ell(\mathbf{s}_y) = \max_{\mathbf{s} \in F_y}\ell(\mathbf{s}).
\end{align*} Since $\pi(\mathbf{x}) = 0$, we have \begin{align*}
    \pi(\mathbf{s}_y + \mathbf{x}) = \pi(\mathbf{s}_y) = y
    \qquad \textcolor{gray}{\left(\text{since } \mathbf{s}_y \in F_y = S \cap \pi^{-1}(y) \subseteq \pi^{-1}(y)\right)}.
\end{align*} Moreover, since $\ell$ is linear and $\ell(\mathbf{x}) > 0$, it follows that \begin{align*}
    \ell(\mathbf{s}_y + \mathbf{x}) = \ell(\mathbf{s}_y) + \ell(\mathbf{x}) > \ell(\mathbf{s}_y).
\end{align*} Hence, by the maximality of $\mathbf{s}_y$ in $F_y$, we have $\mathbf{s}_y + \mathbf{x} \notin S$. That is, \begin{align*}
    \mathbf{s}_y + \mathbf{x} \in (S+\mathbf{x}) \setminus S.
\end{align*} Moreover, if $y_1,y_2 \in \pi(S)$ are such that $y_1 \neq y_2$, then $\pi(\mathbf{s}_{y_1}) = y_1 \neq y_2 = \pi(\mathbf{s}_{y_2})$ and therefore $\mathbf{s}_{y_1} \neq \mathbf{s}_{y_2}$, so that $\mathbf{s}_{y_1} + \mathbf{x} \neq \mathbf{s}_{y_2} + \mathbf{x}$ as well.
And so, \begin{align*}
    y \mapsto \mathbf{s}_y + \mathbf{x}
\end{align*} defines an injection from $\pi(S)$ into $(S+\mathbf{x}) \setminus S$; consequently, \begin{align}\label{injection}
    \abs{(S+\mathbf{x}) \setminus S} \geq \abs{\pi(S)}.
\end{align} 

Next, by linearity of $\pi$, \begin{align*}
    \pi(S) &= \pi\tup{\Sigma(X')} = \set{\pi \tup{\sum_{i \in A}\mathbf{x}_i } : A \subseteq [m-1]} \\ &= \set{\sum_{i \in A}\pi(\mathbf{x}_i) : A \subseteq [m-1]} = \Sigma\tup{\pi(X')}.
\end{align*} We thus obtain \begin{align}\label{induct_X'}
    |\pi(S)| = \abs{\Sigma(\pi(X'))} \geq \abs{\mathcal{I}(\pi(X'))}
\end{align}
by applying the induction hypothesis to $\pi(X')$.
From \eqref{injection} and \eqref{induct_X'}, we obtain \begin{align}\label{ineq_S+x-S}
    \abs{(S+\mathbf{x}) \setminus S} \geq \abs{\pi(S)} \geq \abs{\mathcal{I}(\pi(X'))}.
\end{align} And so, \eqref{split_sigmaX} yields \begin{align}
    \abs{\Sigma(X)} &= \underbrace{\abs{S}}_{\textcolor{gray}{\substack{\geq \abs{\mathcal{I}(X')} \\ \text{(by (\ref{induct_X}))}}}}+ \ \underbrace{\abs{\tup{S + \mathbf{x}} \setminus S}}_{\textcolor{gray}{\substack{\geq \abs{\mathcal{I}(\pi(X'))} \\ \text{(by (\ref{ineq_S+x-S}))}}}} \qquad  \nonumber \\
    &\geq \abs{\mathcal{I}(X')} + \abs{\mathcal{I}(\pi(X'))}.\label{splitI} 
\end{align} 

Finally, let us identify $\mathcal{I}(X)$ in terms of $\mathcal{I}(X')$ and $\mathcal{I}(\pi(X'))$. Every subset $A \subseteq [m]$ either contains $m$ or it does not, and these two cases are disjoint. 

If $m \notin A$, then $A \subseteq [m-1]$; hence, $A \in \mathcal{I}(X)$ if and only if $A \in \mathcal{I}(X')$. 
Thus, the sets $A \in \mathcal{I}\tup{X}$ that don't contain $m$ are exactly the sets $A \in \mathcal{I}\tup{X'}$.

If $m \in A$, let us write $A = B \cup \{m\}$, where $B \subseteq [m-1]$. Then $\tup{\mathbf{x}_i : i \in A}$ is linearly independent if and only if $\tup{\pi(\mathbf{x}_i) : i \in B}$ is linearly independent\footnote{Indeed, if $\sum_{i \in B}c_i\pi(\mathbf{x}_i) = \mathbf{0}$, then linearity of $\pi$ yields $\pi\tup{\sum_{i \in B}c_i\mathbf{x}_i} = \mathbf{0}$. Hence, $\sum_{i \in B}c_i\mathbf{x}_i \in \ker(\pi) = \langle \mathbf{x} \rangle = \langle \mathbf{x}_m \rangle$. Thus, there exists $c_m \in \mathbb{R}$ such that $\sum_{i \in B}c_i\mathbf{x}_i +c_m\mathbf{x}_m = \mathbf{0}$. Conversely, any linear dependence relation among $\tup{\mathbf{x}_i : i \in B \cup \{m\} = A}$ projects under $\pi$ to a linear dependence relation among $\tup{\pi(\mathbf{x}_i) : i \in B}$, since $\mathbf{x}_m = \mathbf{x} \neq \mathbf{0}$ ensures that the $\mathbf{x}_m$-coefficient cannot be the only nonzero coefficient.}, meaning that $B \in \mathcal{I}\tup{\pi(X')}$.
In other words, $A \in \mathcal{I}\tup{X}$ if and only if $B \in \mathcal{I}\tup{\pi(X')}$.
Thus, the sets $A \in \mathcal{I}\tup{X}$ that contain $m$ are in bijection with the sets $B \in \mathcal{I}\tup{\pi(X')}$.

Combining the results of both cases, we find
\begin{align*}
    \mathcal{I}(X) = \underbrace{\mathcal{I}(X')}_{\textcolor{gray}{\text{sets $A$ not containing $m$}}} \sqcup\ \ \ \ \underbrace{\set{B \cup \{m\} : B \in \mathcal{I}(\pi(X'))}}_{\textcolor{gray}{\text{sets $A$ containing $m$}}},
\end{align*} so that
\begin{align*}
    \abs{\mathcal{I}(X)} = \abs{\mathcal{I}(X')} + \abs{\mathcal{I}(\pi(X'))}.
\end{align*} And thus, from (\ref{splitI}), we obtain \begin{align*}
    \abs{\Sigma(X)} \geq \abs{\mathcal{I}(X')} + \abs{\mathcal{I}(\pi(X'))} = \abs{\mathcal{I}(X)}.
\end{align*} This proves (\ref{subset_sums_index_sets}).
\end{proof}

\end{document}
